%  LaTeX support: latex@mdpi.com 
%  For support, please attach all files needed for compiling as well as the log file, and specify your operating system, LaTeX version, and LaTeX editor.

%=================================================================
%\documentclass[journal,article,submit,pdftex,moreauthors]{Definitions/mdpi}
%\documentclass[preprints,article,submit,pdftex,moreauthors]{Definitions/mdpi}
\documentclass[electronics,article,submit,pdftex,moreauthors]{Definitions/mdpi}
% For posting an early version of this manuscript as a preprint, you may use "preprints" as the journal. Changing "submit" to "accept" before posting will remove line numbers.

%=================================================================
% MDPI internal commands - do not modify
\firstpage{1} 
\makeatletter 
\setcounter{page}{\@firstpage} 
\makeatother
\pubvolume{1}
\issuenum{1}
\articlenumber{0}
\pubyear{2026}
\copyrightyear{2026}
%\externaleditor{Firstname Lastname} % More than 1 editor, please add `` and '' before the last editor name
\datereceived{ } 
\daterevised{ } % Comment out if no revised date
\dateaccepted{ } 
\datepublished{ } 
%\datecorrected{} % For corrected papers: "Corrected: XXX" date in the original paper.
%\dateretracted{} % For retracted papers: "Retracted: XXX" date in the original paper.
%\doinum{} % Used for some special journals, like molbank
%\pdfoutput=1 % Uncommented for upload to arXiv.org
%\CorrStatement{yes}  % For updates
%\longauthorlist{yes} % For many authors that exceed the left citation part
%\IsAssociation{yes} % For association journals

%=================================================================
% Add packages and commands here. The following packages are loaded in our class file: fontenc, inputenc, calc, indentfirst, fancyhdr, graphicx, epstopdf, lastpage, ifthen, float, amsmath, amssymb, lineno, setspace, enumitem, mathpazo, booktabs, titlesec, etoolbox, tabto, xcolor, colortbl, soul, multirow, microtype, tikz, totcount, changepage, attrib, upgreek, array, tabularx, pbox, ragged2e, tocloft, marginnote, marginfix, enotez, amsthm, natbib, hyperref, cleveref, scrextend, url, geometry, newfloat, caption, draftwatermark, seqsplit
% cleveref: load \crefname definitions after \begin{document}

\usepackage{rotating}

%\usepackage{adjustbox}
\usepackage[graphicx]{realboxes}
%\usepackage{graphicx}

\usepackage{pdflscape}
\usepackage{afterpage}
\usepackage{booktabs}
%\usepackage[table,xcdraw]{xcolor}
\usepackage{multirow}
\usepackage{makecell}

\usepackage{bm}
%\usepackage{subcaption}

\usepackage{overpic}

\usepackage{comment}

%\usepackage{Definitions/trackchanges}

%\newcommand{\cbf}[1]{\color[HTML]{FE0000} \textbf{#1}}
%\newcommand{\cbf}[1]{\color[HTML]{FE0000} \boldmath{#1}}
\newcommand{\cbfm}[1]{\color[HTML]{FE0000} \boldmath{$#1$}}
\newcommand{\bfm}[1]{\boldmath{$#1$}}

\newcommand{\ccd}{\cellcolor[HTML]{A6A6A6}} % dark (best)
\newcommand{\cc}{\cellcolor[HTML]{D9D9D9}}   % light (second)

\newcommand{\tsup}[1]{\textsuperscript{#1}}
\newcommand{\tsub}[1]{\textsubscript{#1}}



\DeclareUnicodeCharacter{2212}{\textminus}

% -------- Revision highlighting --------
\definecolor{revblue}{RGB}{0,0,210}
\DeclareRobustCommand{\rev}[1]{{\color{revblue}#1}} % (blue color)

%=================================================================
% Please use the following mathematics environments: Theorem, Lemma, Corollary, Proposition, Characterization, Property, Problem, Example, ExamplesandDefinitions, Hypothesis, Remark, Definition, Notation, Assumption
%% For proofs, please use the proof environment (the amsthm package is loaded by the MDPI class).

%=================================================================
% Full title of the paper (Capitalized)
\Title{Generating Diverse Synthetic Eyeglass Datasets Using GANs}


% Author Orchid ID: enter ID or remove command
\newcommand{\orcidauthorA}{0009-0003-3496-9928} % G.S.
\newcommand{\orcidauthorB}{0009-0001-7098-2148} % H.G.
\newcommand{\orcidauthorC}{0000-0001-6669-4445} % R.P.
\newcommand{\orcidauthorD}{0000-0002-5137-4585} % D.M.

% Authors, for the paper (add full first names)
\Author{Giga Shubitidze~*\orcidA{}; Henrikas Giedra~\orcidB{}; Raimondas Pomarnacki~\orcidC{}; and Dalius Matuzevičius~\orcidD{}}

%\longauthorlist{yes}

% MDPI internal command: Authors, for metadata in PDF
\AuthorNames{Giga Shubitidze, Henrikas Giedra, Raimondas Pomarnacki, Dalius Matuzevičius}


% Affiliations / Addresses (Add [1] after \address if there is only one affiliation.)
\address[1]{%
Department of Electronic Systems, Vilnius Gediminas Technical University (VILNIUS TECH), 10105~Vilnius,~Lithuania}

% Contact information of the corresponding author
\corres{Correspondence: author@vilniustech.lt}

% Current address and/or shared authorship
%\firstnote{Current address: Affiliation.}  
% Current address should not be the same as any items in the Affiliation section.
%\secondnote{These authors contributed equally to this work.}
% The commands \thirdnote{} till \eighthnote{} are available for further notes.

%\simplesumm{} % Simple summary

%\conference{} % An extended version of a conference paper

% Abstract (Do not insert blank lines, i.e. \\)
%A single paragraph of about 200 words maximum. For research articles, abstracts should give a pertinent overview of the work. We strongly encourage authors to use the following style of structured abstracts, but without headings:
%(1) Background: place the question addressed in a broad context and highlight the purpose of the study;
%(2) Methods: describe briefly the main methods or treatments applied;
%(3) Results: summarize the article's main findings;
%(4) Conclusions: indicate the main conclusions or interpretations.
%The abstract should be an objective representation of the article, it must not contain results which are not presented and substantiated in the main text and should not exaggerate the main conclusions.

\abstract{[Abstract to be written once the results sections are final.]}



% Keywords
%\keyword{keyword 1; keyword 2; keyword 3 (List three to ten pertinent keywords specific to the article; yet reasonably common within the subject discipline.)}
\keyword{generative adversarial networks; StyleGAN2; dataset diversity; synthetic data generation; generative model evaluation; quality-diversity trade-off; eyeglass frames; limited data}


% The fields PACS, MSC, and JEL may be left empty or commented out if not applicable
%\PACS{J0101}
%\MSC{}
%\JEL{}

%%%%%%%%%%%%%%%%%%%%%%%%%%%%%%%%%%%%%%%%%%
% Only for the journal Applied Sciences
%\featuredapplication{Authors are encouraged to provide a concise description of the specific application or a potential application of the work. This section is not mandatory.}
%%%%%%%%%%%%%%%%%%%%%%%%%%%%%%%%%%%%%%%%%%

%%%%%%%%%%%%%%%%%%%%%%%%%%%%%%%%%%%%%%%%%%
% Only for the journal Data
%\dataset{DOI number or link to the deposited data set if the data set is published separately. If the data set shall be published as a supplement to this paper, this field will be filled by the journal editors. In this case, please submit the data set as a supplement.}
%\datasetlicense{License under which the data set is made available (CC0, CC-BY, CC-BY-SA, CC-BY-NC, etc.)}


%%%%%%%%%%%%%%%%%%%%%%%%%%%%%%%%%%%%%%%%%%
\begin{document}

%%%%%%%%%%%%%%%%%%%%%%%%%%%%%%%%%%%%%%%%%%
%\setcounter{section}{-1} %% Remove this when starting to work on the template.
%\section{How to Use this Template}

%The template details the sections that can be used in a manuscript. Note that the order and names of article sections may differ from the requirements of the journal (e.g., the positioning of the Materials and Methods section). Please check the instructions on the authors' page of the journal to verify the correct order and names. For any questions, please contact the editorial office of the journal or support@mdpi.com. For LaTeX-related questions please contact latex@mdpi.com.%\endnote{This is an endnote.} % To use endnotes, please un-comment \printendnotes below (before References). Only journal Laws uses \footnote.

% The order of the section titles is different for some journals. Please refer to the "Instructions for Authors” on the journal homepage.

\section{Introduction}
\label{sec:intro}

% Synthetic catalogue data for eyeglass frames, and why diversity rather
% than fidelity alone is what makes a generated dataset useful downstream.

\subsection{Motivation}

\subsection{Contributions}

%%%%%%%%%%%%%%%%%%%%%%%%%%%%%%%%%%%%%%%%%%
\section{Related Work}
\label{sec:related}

% StyleGAN2 / StyleGAN2-ADA; GAN evaluation (FID, KID, precision-recall,
% density-coverage, Vendi, LPIPS); training GANs on limited data;
% rejection sampling.

%%%%%%%%%%%%%%%%%%%%%%%%%%%%%%%%%%%%%%%%%%
\section{Materials and Methods}
\label{sec:methods}

\subsection{Dataset}
\label{subsec:data}

The training corpus consists of 6,511 frontal catalogue photographs of
eyeglass frames, the same collection used in our earlier work on joint
image--mask synthesis~\cite{shubitidze2026masks}. Each source image is square,
between roughly 1350 and 1650 pixels on a side with a median of 1463, and
shows a single unworn frame, roughly centred, against a near-uniform light
background. Images are centre-cropped
to a 2:1 aspect ratio and resized to $256 \times 512$ pixels (height $\times$
width), which follows the natural proportions of a frame viewed from the
front and avoids the padding a square canvas would demand.

Two properties of this corpus shape everything that follows. The first is the
background. Because it is close to uniform, a frame silhouette can be
recovered by thresholding on distance from the background colour, with no
learned segmentation and without asking the generator to emit a mask. Every
geometric and symmetry measure in Section~\ref{sec:metrics} is built on
silhouettes obtained this way, which is what makes a geometry-aware diversity
study possible in a setting where mask generation is deliberately out of
scope. The second property is size. At 6,511 images the corpus is small for
adversarial training at this resolution, and in practice the majority of the
failures reported below trace back to the discriminator memorising it rather
than to any limitation of the generator.

Two dataset-level interventions are used throughout. Horizontal flips are
applied to the dataset rather than as a stochastic augmentation, doubling the
effective corpus to roughly 13,000 images. Sampling is additionally weighted
by inverse frequency over a $16 \times 16$ histogram of each frame's mean
CIELAB $a^*b^*$ chromaticity, with exponent 0.5 and a weight cap of 4. The
diagnosis that motivated this weighting, and the dose--response behaviour it
produces, are treated in Section~\ref{subsec:ablations}; it is introduced here
because every model from that point onward was trained under it.

\subsection{Base Model and Stabilisation}
\label{subsec:base}

The starting point is the conditional StyleGAN2 variant of our previous
work~\cite{shubitidze2026masks}, which generated an RGB image and a binary
frame mask from a single forward pass and reached an FID of 6.13. A later
revision of that model added an eight-layer mapping network, adaptive
discriminator augmentation~\cite{karras2020ada}, path-length regularisation, a
feature-matching term and KID-based checkpoint selection. It produced visibly
sharper samples and the best domain metrics of that work, but its
distributional and diversity behaviour was worse rather than better: FID rose
to 38.47, and the colour-diversity score fell to 0.0567 against 0.0995 for the
earlier run.

That comparison also indicates why this paper does not simply reuse the
earlier diversity metrics. The colour-diversity measure of the previous work
is a mean pairwise distance computed over generated images alone, and on the
real corpus it evaluates to 0.0338. Both generators therefore scored well
above the real data on a measure intended to certify their realism, which a
statistic of that construction permits: it rewards spread without reference to
where the real distribution lies. A quantity that can be improved without
limit by drifting away from the target does not measure how well the target is
covered, and Section~\ref{sec:metrics} replaces it.

Removing the mask branch, as the present scope requires, destabilised training
badly, and three defects had to be corrected before a stable unconditional
baseline existed. The first was the output head. The generator projected to
RGB through a $\tanh$ nonlinearity, inherited from the mask-supervised model.
Approximately 86 per cent of the pixels in a cropped and resized catalogue
image are near-white background, so every loss term drives the large majority
of pixels toward $+1$, where $\tanh$ has vanishing gradient. The official StyleGAN2 generator uses a linear RGB
head~\cite{karras2020stylegan2}, and adopting it removed the saturation.

The second defect was in the mapping network, and it is the one with a
plausible bearing on the low colour diversity of the predecessor run. Eight
stacked default linear layers with LeakyReLU activations contract the signal
at every step, because standard initialisation scales weights by the inverse
square root of the fan-in without the equalized-learning-rate correction.
Measured at initialisation on a batch of 4096 standard-normal latents, the
resulting style codes have a standard deviation of 0.0184 against 1.0000 for
the latents themselves, and the contraction is front-loaded: the per-layer
standard deviation falls from 0.583 after the first linear layer to 0.052
after the fourth before levelling off near 0.028. A style code of that
dispersion carries almost no per-image information, which leaves per-pixel
noise injection as the dominant source of variation between samples. Replacing
the stack with equalized-learning-rate linear layers at a multiplier of 0.01,
as in the official mapping network~\cite{karras2020stylegan2}, raises the
style-code standard deviation to 1.0968 under the identical measurement.

The third defect was the adversarial loss. The hinge
formulation~\cite{lim2017geometric} produces no gradient once real and fake
scores separate beyond its margins. Under mask supervision the discriminator
had never built margins wide enough for this to matter; without it, the two
distributions separate within a few epochs, the discriminator loss reaches
exactly zero, and training stalls. Substituting the non-saturating logistic
loss standard in StyleGAN2-ADA~\cite{karras2020stylegan2,karras2020ada}
removed the stall.

\subsection{Training Configuration}
\label{subsec:training}

All models reported here are StyleGAN2-ADA derivatives implemented in stock
PyTorch~\cite{paszke2019pytorch}, with no compiled custom operators. The
generator maps a 512-dimensional latent through the eight-layer
equalized-learning-rate mapping network described above to a 512-dimensional
style code, and synthesises through six weight-demodulated stages from $8
\times 16$ to $256 \times 512$ with skip-RGB connections and a linear output
head. A self-attention module~\cite{zhang2019self} is placed at $32 \times 64$
to give the network access to long-range structure, which for frontal frames
is chiefly bilateral correspondence. Channel widths follow $512 \rightarrow
256 \rightarrow 256 \rightarrow 128 \rightarrow 64 \rightarrow 32 \rightarrow
16$, giving 8.50~M generator parameters; the discriminator uses $16
\rightarrow 32 \rightarrow 64 \rightarrow 128 \rightarrow 256 \rightarrow 256$
for 2.50~M parameters. Both width multipliers were varied, and the outcome is
reported in Section~\ref{subsec:negative}.

Training uses the non-saturating logistic loss with a lazy $R_1$ gradient
penalty~\cite{mescheder2018converge} at $\gamma = 5$ applied every 16 steps,
path-length regularisation at weight 4.0 on the same interval with a reduced
batch of 16 to fit memory, a feature-matching term at weight 2.0 with random
real--fake pairing, and a VGG perceptual term~\cite{simonyan2015vgg} at weight
0.75 with nearest-real pairing. The discriminator applies a group-wise
minibatch standard deviation layer~\cite{karras2018progressive} over groups of
four and blurs feature maps before average pooling. Adaptive discriminator
augmentation uses horizontal flips, integer-pixel translation, brightness
jitter, additive noise and cutout, with the probability ceiling at 0.95; hue
and saturation augmentation is disabled, for reasons established in
Section~\ref{subsec:negative}. Both networks use
Adam~\cite{kingma2015adam} with $\beta_1 = 0$ and $\beta_2 = 0.99$ at learning
rates of $2 \times 10^{-4}$ for the generator and $5 \times 10^{-5}$ for the
discriminator, style mixing is applied with probability 0.9, and an
exponential moving average of the generator weights with decay 0.999 is used
for all sampling and evaluation. The batch size is 16. On a single RTX~4090
with 24~GB of memory an epoch takes approximately 180 to 190 seconds.

\subsection{Generation Protocol}
\label{subsec:generation}

Sampling is a design space in its own right, independent of training, and it
was swept separately. Five levers were considered: the truncation coefficient
$\psi$; layered truncation, in which coarse and fine synthesis stages receive
different coefficients; whether each image is synthesised from a single latent
or from a blend of several style codes; multi-modal truncation about $k$
cluster centres rather than the global style mean; and post-hoc rejection
sampling against either a manifold-precision criterion or the structural
artifact suite of Section~\ref{subsec:artifacts}. Nine protocols built from
these levers were evaluated on a common checkpoint
(Table~\ref{tab:protocols}), each generating 10,000 images at a fixed seed.

\begin{table}[H]
\caption{Generation protocols evaluated on a common checkpoint. Yield is the
fraction of the 10,000 generated images surviving rejection sampling.
Truncation is applied toward the style mean; ``pure'' denotes a single latent
per image and ``mixed'' the blended default; ``layered'' applies $\psi = 1.0$
to the three coarse stages and $\psi = 1.2$ to the three fine stages.}
\label{tab:protocols}
\setlength{\tabcolsep}{5pt}
\begin{tabular}{lcccccc}
\toprule
\textbf{Protocol} & \textbf{FID} $\downarrow$ & \textbf{Prec.} &
\textbf{Recall} & \textbf{$a^*b^*$ cov.} & \textbf{Vendi} & \textbf{Yield} \\
\midrule
$\psi = 0.85$, mixed                 & 13.09 & 0.963 & 0.103 & 0.219 & 2.29 & 100\% \\
$\psi = 1.0$, mixed                  & 11.25 & 0.924 & 0.154 & 0.270 & 2.52 & 100\% \\
$\psi = 1.2$, mixed                  & 10.46 & 0.844 & 0.235 & 0.352 & 2.83 & 100\% \\
$\psi = 1.0$, pure                   & 10.41 & 0.890 & 0.189 & 0.313 & 2.70 & 100\% \\
$\psi = 1.2$, pure                   & 10.76 & 0.765 & 0.324 & 0.395 & 3.05 & 100\% \\
$\psi = 1.2$, pure, precision filter & 10.35 & 0.923 & 0.250 & 0.355 & 2.85 & 69.7\% \\
Layered $1.0/1.2$, pure              & 10.20 & 0.879 & 0.212 & 0.367 & 2.74 & 100\% \\
Layered, precision filter            & 10.51 & 0.957 & 0.169 & 0.344 & 2.65 & 82.9\% \\
Layered, $k = 8$ centres             & 10.22 & 0.879 & 0.203 & 0.352 & 2.73 & 100\% \\
\bottomrule
\end{tabular}
\end{table}

Three findings follow from this sweep, and one of them runs against common
practice. Truncation below $\psi = 1$ is not a trade-off in this domain but a
straightforward loss: moving from 0.85 to 1.0 to 1.2 improves FID and every
diversity measure at once, so the low-$\psi$ settings conventionally used to
buy image quality~\cite{brock2019biggan} do not lie on the front at all.
Synthesising each image from a single latent outperforms blending several
style codes at matched $\psi$ on nearly every axis, which is unsurprising once
stated, since averaging style codes moves each sample toward the population
mean, but the blended scheme was the default inherited from our previous work.
Layered truncation gives the best compromise: capping only the three coarse
stages, which carry silhouette and proportion, while leaving untruncated the
three fine stages that carry colour and texture yields the lowest FID at full
yield, and it beats the filtered $\psi = 1.2$ set on colour coverage precisely
because the stages responsible for colour are never contracted.

Two levers did not repay their cost. Multi-modal truncation about eight
centres differs from the single-centre layered protocol by 0.02 in FID, which
is inside the measurement noise established in Section~\ref{subsec:noise}; at
$\psi \geq 1$ the contraction is evidently too weak for the choice of centre
to matter. Rejection sampling, useful earlier, became unnecessary once layered
truncation had removed the manifold-outlier tail, and stacking it on top moves
along a trade-off curve rather than lifting it; it is examined on its own
terms in Section~\ref{subsec:filtering}.

The layered protocol was frozen as the headline configuration and used for
every model comparison in this paper: $\psi = 1.0$ for the coarse stages and
$\psi = 1.2$ for the fine stages with the split after the third stage, a
single latent per image, 10,000 images, generation seed 42. One caveat must be
stated plainly. The protocol was selected on a checkpoint whose FID was 10.2,
whereas the final models reach 3.6 to 4.1; it has not been re-validated at
that operating point, and the model by protocol coverage discussed in
Section~\ref{subsec:matrix} is too sparse to establish that the ranking
transfers. Holding the protocol fixed makes the model comparison internally
consistent, which is what that comparison requires, but it does not
demonstrate that this protocol is optimal for the final models.

\subsection{Reproducibility}
\label{subsec:reproducibility}

Bitwise reproduction of these runs is not achievable and is not claimed.
Neither cuDNN autotuning nor non-deterministic reduction kernels are pinned,
and because each resume re-seeds the random number generator from the
configured seed rather than restoring generator state, the batch ordering
after an interruption is not what an uninterrupted run would have produced.
The claim we make is at the level of method: a fixed seed, a pinned commit,
and a published configuration and dataset. Run-to-run variation under a
different training seed was never measured. Quantifying it would require
repeating a full run at roughly 63 GPU-hours, and we record it as a limitation
rather than estimate it.

Provenance is recorded automatically. From the point at which run stamping was
introduced, every configuration snapshot, diversity report and artifact report
carries the git commit, a flag indicating an unclean working tree, and the
PyTorch, CUDA and GPU version strings, and every launch appends its full
argument list to a per-run history file. Runs predating that change have
reconstructed provenance only. The longest run illustrates why this matters:
what is reported here as a single 1,600-epoch training was in fact executed as
six consecutive segments, and its structure had to be recovered after the fact
from resets in the cumulative training-hours series, giving segments of 300,
200, 300, 357, 43 and 400 epochs for a total of 83.9 hours. Epochs are
reported throughout as the number of training epochs completed; checkpoint
files store a zero-based index one lower, so the checkpoint behind epoch 1600
records 1599.

%%%%%%%%%%%%%%%%%%%%%%%%%%%%%%%%%%%%%%%%%%
\section{A Diversity Metric Suite for Eyeglass Frames}
\label{sec:metrics}

% Core contribution. Kept as a top-level section rather than folded into
% Methods, since the two new metrics are the point of the paper.

\subsection{Geometry, Colour, Feature-Space and Perceptual Measures}
\label{subsec:metric_families}

\subsection{Pose-Aligned Mirror Symmetry}
\label{subsec:aligned_sym}

% Include the synthetic validation table: a 2 deg tilt costs 0.51 of the
% fixed-axis score while genuine shape asymmetry costs only 0.07.

\subsection{Mirror-Axis Offset Dispersion}
\label{subsec:yaw}

% Measures differential temple foreshortening (yaw). Translation-invariant
% by construction - state this explicitly, it is easy to misread.

\subsection{Structural Artifact Suite}
\label{subsec:artifacts}

\subsection{Measurement Noise Floor}
\label{subsec:noise}

% Per-metric maximum across the two generation-seed pairs: FID 0.052,
% recall 0.019, ab_coverage 0.016, precision 0.008, coverage 0.003.

%%%%%%%%%%%%%%%%%%%%%%%%%%%%%%%%%%%%%%%%%%
\section{Results}
\label{sec:results}

\subsection{Training Ablations}
\label{subsec:ablations}

% Organised by mechanism, not chronologically: R1 gamma, colour rebalancing
% dose-response, discriminator blurpool, mirror coupling.

\subsection{Quality-Diversity Pareto Front}
\label{subsec:pareto}

% Two headline models, neither dominating the other - that is the finding.
% Figure: images/pareto.pdf. Table: results_table.tex.

\subsection{Model and Generation-Strategy Coverage}
\label{subsec:matrix}

% Cut from the paper. The coverage is too sparse and too non-factorial to
% support a model x strategy claim; the caveat is stated in Section 3.4.

\subsection{Rejection Filtering as a Pareto Knob}
\label{subsec:filtering}

% Filtering traces a front between structural cleanliness and colour
% coverage, since defects concentrate in the rare-colour tail.

%%%%%%%%%%%%%%%%%%%%%%%%%%%%%%%%%%%%%%%%%%
\section{Checkpoint Selection Under a Diversity Objective}
\label{sec:checkpoints}

% Self-contained study: does the minimum-KID checkpoint give the best
% dataset? Four groups over 16 usable runs (5/4/2/5), plus the epoch curve.

%%%%%%%%%%%%%%%%%%%%%%%%%%%%%%%%%%%%%%%%%%
\section{Discussion}
\label{sec:discussion}

\subsection{Preprocessing Mismatch and Measured Quality}
\label{subsec:resize}

% The strongest methodological finding: a train/eval resize filter mismatch
% capped measured output quality across six architecture experiments and
% produced three plausible but wrong explanations before it was located.

\subsection{Silent Failure Modes in Training and Checkpointing}
\label{subsec:silent}

% Silent AMP generator-step skipping; partial EMA load; and a checkpoint
% invalidated by a code change with no key or shape mismatch.

\subsection{Negative Results}
\label{subsec:negative}

% Falsified hypotheses: weaker D, colour ADA, D input channel, G capacity,
% and the yaw reading behind the second mirror variant.

\subsection{Method Lessons}
\label{subsec:lessons}

% What we would do differently: validate a new metric on controlled input
% before building on it; change one thing at a time; diff every launch.

\subsection{Limitations}
\label{subsec:limitations}

% Reproducibility limits, and yaw collapse under a symmetry prior as a
% measured property of the prior rather than a defect of one model.

%%%%%%%%%%%%%%%%%%%%%%%%%%%%%%%%%%%%%%%%%%
\section{Conclusions}
\label{sec:conclusions}

%%%%%%%%%%%%%%%%%%%%%%%%%%%%%%%%%%%%%%%%%%
%\section{Patents}
%
%This section is not mandatory, but may be added if there are patents resulting from the work reported in this manuscript.

%%%%%%%%%%%%%%%%%%%%%%%%%%%%%%%%%%%%%%%%%%
\vspace{6pt} 

%%%%%%%%%%%%%%%%%%%%%%%%%%%%%%%%%%%%%%%%%%
%% optional
%\supplementary{The following supporting information can be downloaded at:  \linksupplementary{s1}, Figure S1: title; Table S1: title; Video S1: title.}


%%%%%%%%%%%%%%%%%%%%%%%%%%%%%%%%%%%%%%%%%%
\authorcontributions{Conceptualization, G.S. and D.M.; methodology, G.S. and D.M.; software, G.S.; validation, G.S., H.G., R.P. and D.M.; formal analysis, G.S.; investigation, G.S.; resources, D.M.; data curation, G.S.; writing---original draft preparation, G.S.; writing---review and editing, G.S., H.G., R.P. and D.M.; visualization, G.S.; supervision, D.M.; project administration, D.M. All authors have read and agreed to the published version of the manuscript.}


%\funding{Please add: ``This research received no external funding'' or ``This research was funded by NAME OF FUNDER grant number XXX.'' and  and ``The APC was funded by XXX''. Check carefully that the details given are accurate and use the standard spelling of funding agency names at \url{https://search.crossref.org/funding}, any errors may affect your future funding.}
% Grant carried over from the previous article - confirm it still applies.
\funding{This research was supported by the Research Council of Lithuania (LMTLT), agreement No S-ITP-24-12.}

%\institutionalreview{In this section, you should add the Institutional Review Board Statement and approval number, if relevant to your study. You might choose to exclude this statement if the study did not require ethical approval. Please note that the Editorial Office might ask you for further information. Please add “The study was conducted in accordance with the Declaration of Helsinki, and approved by the Institutional Review Board (or Ethics Committee) of NAME OF INSTITUTE (protocol code XXX and date of approval).” for studies involving humans. OR “The animal study protocol was approved by the Institutional Review Board (or Ethics Committee) of NAME OF INSTITUTE (protocol code XXX and date of approval).” for studies involving animals. OR “Ethical review and approval were waived for this study due to REASON (please provide a detailed justification).” OR “Not applicable” for studies not involving humans or animals.}
\institutionalreview{Not applicable.}

%\informedconsent{Any research article describing a study involving humans should contain this statement. Please add ``Informed consent was obtained from all subjects involved in the study.'' OR ``Patient consent was waived due to REASON (please provide a detailed justification).'' OR ``Not applicable'' for studies not involving humans. You might also choose to exclude this statement if the study did not involve humans.
\informedconsent{Not applicable.}
%
%Written informed consent for publication must be obtained from participating patients who can be identified (including by the patients themselves). Please state ``Written informed consent has been obtained from the patient(s) to publish this paper'' if applicable.}

%\dataavailability{We encourage all authors of articles published in MDPI journals to share their research data. In this section, please provide details regarding where data supporting reported results can be found, including links to publicly archived datasets analyzed or generated during the study. Where no new data were created, or where data is unavailable due to privacy or ethical restrictions, a statement is still required. Suggested Data Availability Statements are available in section ``MDPI Research Data Policies'' at \url{https://www.mdpi.com/ethics}.} 
%\dataavailability{Image datasets used in this research are described in Section~\ref{ssec:2MM:data_preparation}.}
%\dataavailability{All data underlying the results are available as part of the article and no additional source data are required.}
%\dataavailability{The data that support the findings of this study are available from the corresponding author, H.G., upon reasonable request.}
\dataavailability{The raw data supporting the conclusions of this article will be made available by the corresponding author on request.}

% Proposed during submission: The data presented in this study are openly available in [repository name, e.g., FigShare] at [doi], reference number [reference number].
%\acknowledgments{In this section you can acknowledge any support given which is not covered by the author contribution or funding sections. This may include administrative and technical support, or donations in kind (e.g., materials used for experiments).}


%\conflictsofinterest{Declare conflicts of interest or state ``The authors declare no conflicts of interest.'' Authors must identify and declare any personal circumstances or interest that may be perceived as inappropriately influencing the representation or interpretation of reported research results. Any role of the funders in the design of the study; in the collection, analyses or interpretation of data; in the writing of the manuscript; or in the decision to publish the results must be declared in this section. If there is no role, please state ``The funders had no role in the design of the study; in the collection, analyses, or interpretation of data; in the writing of the manuscript; or in the decision to publish the results''.} 
\conflictsofinterest{The authors declare no conflicts of~interest.}

%%%%%%%%%%%%%%%%%%%%%%%%%%%%%%%%%%%%%%%%%%
\abbreviations{Abbreviations}{
The following abbreviations are used in this manuscript:\\

\noindent 
\begin{tabular}{@{}ll}
ADA & Adaptive discriminator augmentation\\
CIELAB & CIE $L^*a^*b^*$ colour space\\
EMA & Exponential moving average\\
FID & Fr\'echet inception distance\\
GAN & Generative adversarial network\\
KID & Kernel inception distance\\
LPIPS & Learned perceptual image patch similarity\\
PPL & Perceptual path length\\
PRDC & Precision, recall, density and coverage\\
RGB & Red, green, blue\\
VRAM & Video random-access memory\\
\end{tabular}
}

%%%%%%%%%%%%%%%%%%%%%%%%%%%%%%%%%%%%%%%%%%
\appendixtitles{no}
%\appendixstart
%\appendix
%\section{Appendix Title}
%\label{app:example}

%\finishlandscape

%%%%%%%%%%%%%%%%%%%%%%%%%%%%%%%%%%%%%%%%%%
\begin{adjustwidth}{-\extralength}{0cm}
%\printendnotes[custom] % Un-comment to print a list of endnotes

\reftitle{References}

% Please provide either the correct journal abbreviation (e.g. according to the “List of Title Word Abbreviations” http://www.issn.org/services/online-services/access-to-the-ltwa/) or the full name of the journal.
% Citations and References in Supplementary files are permitted provided that they also appear in the reference list here. 

%=====================================
% References, variant A: external bibliography
%=====================================
\bibliography{refs}

%=====================================
% References, variant B: internal bibliography
%=====================================
% egin{thebibliography}{999}
% ibitem[Author(Year)]{key} ...
% \end{thebibliography}




% If authors have biography, please use the format below
%\section*{Short Biography of Authors}
%\bio
%{\raisebox{-0.35cm}{\includegraphics[width=3.5cm,height=5.3cm,clip,keepaspectratio]{Definitions/author1.pdf}}}
%{\textbf{Firstname Lastname} Biography of first author}
%
%\bio
%{\raisebox{-0.35cm}{\includegraphics[width=3.5cm,height=5.3cm,clip,keepaspectratio]{Definitions/author2.jpg}}}
%{\textbf{Firstname Lastname} Biography of second author}

% For the MDPI journals use author-date citation, please follow the formatting guidelines on http://www.mdpi.com/authors/references
% To cite two works by the same author: \citeauthor{ref-journal-1a} (\citeyear{ref-journal-1a}, \citeyear{ref-journal-1b}). This produces: Whittaker (1967, 1975)
% To cite two works by the same author with specific pages: \citeauthor{ref-journal-3a} (\citeyear{ref-journal-3a}, p. 328; \citeyear{ref-journal-3b}, p.475). This produces: Wong (1999, p. 328; 2000, p. 475)


%%%%%%%%%%%%%%%%%%%%%%%%%%%%%%%%%%%%%%%%%%
\PublishersNote{}
\end{adjustwidth}
\end{document}